\chapter{Introduction}

Out of all the monumental achievements in quantum computing and the different algorithms that have been designed throughout the years, Shor's algorithm perhaps remains as the greatest demonstration of what quantum computers are capable of. Often designated as the crown jewel of the field, it is the first algorithm that truly presents an idea that manages to affect our everyday lives. Its mere existence,along with the possibilities that said existence implies,has even inspired the development of new fields like post-quantum cryptography, that aims to create new challenges upon which modern encryption can be based. All of that purely based on the fact that it completely redefines one of oldest problems in number theory: prime factorization

\chapter{The idea}

First, it is important to define the objective. Given a (usually large) composite number $N$ whose factors are $p$ and $q$ (meaning $N = p\cdot q$), the task is to find the numbers $p$ and $q$. This can in turn be extended to find the full prime factorization of $N$ by recursively applying the algorithm to $p$ and $q$. Hence, we will only be dealing with the case in which $p$ and $q$ are prime numbers. From now and for the entirety of the paper, we define $n = \lceil log_2 N \rceil$ as the number of bits in N.

Furthermore, there are certain trivial cases for $N$. For example, we can easily check whether $N$ is even in $O(logN)$ using Euclid's algorithm, or a prime power ($N = p^k$ for some $p$), with other classical algorithms. This means that we only need to look at the case where $N$ is an odd composite number that is not a prime power.

The classical part of Shor's algorithm begins by picking a random number $a$ with $2<a<N$, and running Euclid's algorithm to compute $gcd(a,N)$. If $\gcd(a,N) > 1$, then we have a non trivial factor,and the other factor can be found by simply dividing $N/\gcd(a,N)$. Otherwise,it follows that $a$ and $N$ are coprime and $a$ is an element of the multiplicative group of the integers modulo N, meaning it has a multiplicative inverse modulo N. This also implies:

\begin{equation}
a^r \equiv 1 \pmod N
\end{equation}
for some $r$ with $1<r<\Phi(N)$\footnote{$r$ is bounded up by $\Phi(N)$ due to Lagrange's theorem}, where r is called the multiplicative order, or period, of $a$ modulo $N$ and $\Phi(N)$ is Euler's totient function. It should be noted that $r$ is defined to be the smallest possible number that satisfies \textbf{Eq 2.1}.

In this case, Shor's algorithm applies a quantum subroutine to find $r$. We will analyze this subroutine in the following chapter, as it directly ties into Regev's algorithm.

From \textbf{Eq. 2.1}, it is obvious that  $N\mid(a^r - 1)$. This implies (but is not equivalent to):

\begin{equation}
N\mid(a^{r/2} - 1)(a^{r/2} + 1)
\end{equation}

From \textbf{Eq 2.2}, a few things can be seen. First,our logic from now on will only work for even $r$. In the case that $r$ happens to be odd, Shor picks a new $a$, computes a new $r$,and relies on probabilistically finding an even $r$.


Also, it cannot be the case that $N\mid(a^{r/2} - 1)$, since that would imply that $a^{r/2} \equiv 1\: \pmod N$, which contradicts the assumption that $r$ is the smallest possible number that satisfies \textbf{Eq 2.1}. At this point,we have two possibilities:either $N\mid(a^{r/2} + 1)\implies a^{r/2} \equiv -1 \pmod N$ ,in which case $\gcd(a^{r/2} -1,N) = 1$, and we don't have any nontrivial factors, or  $N\nmid(a^{r/2} + 1)\implies a^{r/2} \not\equiv -1 \pmod N$ ,and therefore $d = gcd(a^{r/2} -1,N)$ is a nontrivial prime factor of $N$,with the other being $N/d$, and the algorithm is finished.

To summarize the steps we just described:\\
1.Pick a number $a$ between 2 and $N$\\
2.If $gcd(a,N) = 1$,then return $a$ and $N/a$ as nontrivial factors, else run a quantum subroutine to find the order $r$ of $a$ modulo $N$\\
3.If $r$ is even, return to step 1, else compute $d = \gcd(a^{r/2} -1,N)$\\
4.If $d > 1$, return $d$ and $N/d$, otherwise return to step 1.\\










